\documentclass[10pt,a4paper]{amsart}
\usepackage{amsmath,amssymb,amsthm,graphicx}
\usepackage{verbatim}
\newenvironment{reference}{\comment}{\endcomment}
%\newenvironment{reference}{}{}
\newenvironment{slogan}{\comment}{\endcomment}
\newenvironment{history}{\comment}{\endcomment}

% For commutative diagrams we use Xy-pic
\usepackage[all]{xy}

% We use 2cell for 2-commutative diagrams.
\xyoption{2cell}
% Theorem environments.
%
\theoremstyle{plain}
\newtheorem{theorem}[subsection]{Theorem}
\newtheorem{proposition}[subsection]{Proposition}
\newtheorem{lemma}[subsection]{Lemma}

\theoremstyle{definition}
\newtheorem{definition}[subsection]{Definition}
\newtheorem{example}[subsection]{Example}
\newtheorem{exercise}[subsection]{Exercise}
\newtheorem{situation}[subsection]{Situation}

\theoremstyle{remark}
\newtheorem{remark}[subsection]{Remark}
\newtheorem{remarks}[subsection]{Remarks}

\numberwithin{equation}{subsection}

% Macros
%
\def\lim{\mathop{\mathrm{lim}}\nolimits}
\def\colim{\mathop{\mathrm{colim}}\nolimits}
\def\Spec{\mathop{\mathrm{Spec}}}
\def\Hom{\mathop{\mathrm{Hom}}\nolimits}
\def\Ext{\mathop{\mathrm{Ext}}\nolimits}
\def\SheafHom{\mathop{\mathcal{H}\!\mathit{om}}\nolimits}
\def\SheafExt{\mathop{\mathcal{E}\!\mathit{xt}}\nolimits}
\def\Sch{\mathit{Sch}}
\def\Mor{\mathop{\mathrm{Mor}}\nolimits}
\def\Ob{\mathop{\mathrm{Ob}}\nolimits}
\def\Sh{\mathop{\mathit{Sh}}\nolimits}
\def\NL{\mathop{N\!L}\nolimits}
\def\CH{\mathop{\mathrm{CH}}\nolimits}
\def\proetale{{pro\text{-}\acute{e}tale}}
\def\etale{{\acute{e}tale}}
\def\QCoh{\mathit{QCoh}}
\def\Ker{\mathop{\mathrm{Ker}}}
\def\Im{\mathop{\mathrm{Im}}}
\def\Coker{\mathop{\mathrm{Coker}}}
\def\Coim{\mathop{\mathrm{Coim}}}

% Boxtimes
%
\DeclareMathSymbol{\boxtimes}{\mathbin}{AMSa}{"02}

%
% Macros for moduli stacks/spaces
%
\def\QCohstack{\mathcal{QC}\!\mathit{oh}}
\def\Cohstack{\mathcal{C}\!\mathit{oh}}
\def\Spacesstack{\mathcal{S}\!\mathit{paces}}
\def\Quotfunctor{\mathrm{Quot}}
\def\Hilbfunctor{\mathrm{Hilb}}
\def\Curvesstack{\mathcal{C}\!\mathit{urves}}
\def\Polarizedstack{\mathcal{P}\!\mathit{olarized}}
\def\Complexesstack{\mathcal{C}\!\mathit{omplexes}}
% \Pic is the operator that assigns to X its picard group, usage \Pic(X)
% \Picardstack_{X/B} denotes the Picard stack of X over B
% \Picardfunctor_{X/B} denotes the Picard functor of X over B
\def\Pic{\mathop{\mathrm{Pic}}\nolimits}
\def\Picardstack{\mathcal{P}\!\mathit{ic}}
\def\Picardfunctor{\mathrm{Pic}}
\def\Deformationcategory{\mathcal{D}\!\mathit{ef}}
\usepackage{cleveref}
\crefname{originallemma}{original lemma}{original lemmas}
\Crefname{originallemma}{Original lemma}{Original lemmas}
\makeatletter
\newlabel{lemma-essentially-constant}{{070C}{source}{}{sourcepointer-1}{}}
\newlabel{lemma-essentially-constant@cref}{{[originallemma][0][]070C}{[1][0][]source}{}{}{}}
\makeatother
\title{Mittag-Leffler along a short exact sequence of inverse systems with essentially constant quotient}
\author{The Stacks Project Authors}
\date{stacks-project:homology.tex:local-snapshot-20261012; GNU FDL 1.2 or later}
\begin{document}
\maketitle
\noindent\textit{Editorial note.} All delivered mathematical text is the author's own, copied from the pinned Stacks Project chapter file homology.tex; no statement or proof step was rewritten, reordered or removed.\par\medskip
\noindent\textit{Editorial note.} The proof is the author's complete argument; the parenthetical '(small detail omitted)' at the end of the first case closes the routine translation of the stabilisation statement into the Mittag-Leffler condition, whose content is displayed immediately before it. The corpus counts comparable proofs: the card STK-algebra-adjoint-tensor counts sources/stacks-algebra lines 2695-2702, which contain the sentence 'We omit the verification that these constructions are mutually inverse.', and STK-algebra-loc-and-modules counts lines 1345-1353 with the same kind of closing verification omitted. The parenthetical is preserved verbatim inside the counted span.\par\medskip
\noindent\textit{Editorial note.} Every labelled reference occurring inside the counted spans is carried in assembly.external\_references as a non-load-bearing navigation pointer (same-chapter ones at their definition site, cross-chapter ones at the invocation inside the counted span), so no load-bearing mathematics is replaced by a stub. No bibliography citation occurs inside the counted spans.\par\medskip
\noindent\textit{Editorial note.} The same-name check for this card scanned every staged Stacks chapter of this repository -- all sources/stacks-*/pinned-original-source.tex files, the card's own chapter excluded -- and the labels found are recorded in dedup.same\_name\_cross\_chapter on the card.\par\medskip
\noindent\textit{Editorial note.} The preamble is the chapter's own native preamble excerpt already used for this source: the comment-type reference, slogan and history environments (native preamble lines 8-13), the Xy-pic setup (native preamble lines 14-18) and the native theorem declarations with the mathematical macros (native preamble lines 154-217, the actual last line of the file).\par\medskip
\noindent\textit{Allowed standard prerequisites (editorial).}\begin{itemize}
\item Essentially constant inverse systems. Same chapter, Lemma lemma-essentially-constant. Supplies the decomposition of the quotient system as a constant part plus a part whose transition maps are eventually zero, which is the form in which the hypothesis is used.
\item The Mittag-Leffler condition. Same chapter, the definition of ML for an inverse system of abelian groups. Supplies the condition on the stabilisation of the images of the transition maps that is compared for the two systems.
\end{itemize}\medskip
\ifdefined\hypertarget\hypertarget{sourcepointer-1}{}\fi
\noindent\textit{Source reference.} Original lemma 070C: sources/stacks-homology/pinned-original-source.tex, source lines 7950--7950. Navigation pointer invoked at proof line 7991; the counted span contains the pointer only and the cited result is not reproduced inside it, so the pointer is registered as a non-load-bearing navigation stub in the established form for Stacks material. No bibliography citation is involved.\par\medskip
\begin{lemma}
\label{lemma-exact-sequence-ML}
Let
$$
0 \to (A_i) \to (B_i) \to (C_i) \to 0
$$
be an exact sequence of inverse systems of abelian groups.
If $(C_i)$ is essentially constant, then $(A_i)$ has ML
if and only if $(B_i)$ has ML.
\end{lemma}

\begin{proof}
After renumbering we may assume that $C_i = C \oplus Z_i$ compatible with
transition maps and that for all $i$ there exists an $i' \geq i$ such that
$Z_{i'} \to Z_i$ is zero, see Lemma \ref{lemma-essentially-constant}.

\medskip\noindent
First, assume $C = 0$, i.e., we have $C_i = Z_i$.
In this case choose $1 = n_1 < n_2 < n_3 < \ldots$
such that $Z_{n_{i + 1}} \to Z_{n_i}$ is zero.
Then $B_{n_{i + 1}} \to B_{n_i}$ factors through $A_{n_i} \subset B_{n_i}$.
It follows that for $j \geq i + 1$ we have
$$
\Im(A_{n_j} \to A_{n_i}) \subset \Im(B_{n_j} \to B_{n_i}) \subset
\Im(A_{n_{j - 1}} \to A_{n_i})
$$
as subsets of $A_{n_i}$.
Thus the images $\Im(A_{n_j} \to A_{n_i})$ stabilize for $j \geq i + 1$
if and only if the same is true for the images $\Im(B_{n_j} \to B_{n_i})$.
The equivalence follows from this (small detail omitted).

\medskip\noindent
If $C \not = 0$, denote $B'_i \subset B_i$ the inverse image
of $C$ by the map $B_i \to C \oplus Z_i$. Then by the previous
paragraph we see that $(B'_i)$ has ML if and only if $(B_i)$
has ML. Thus we may replace $(B_i)$ by $(B'_i)$.
In this case we have exact sequences $0 \to A_i \to B_i \to C \to 0$
for all $i$. It follows that
$0 \to \Im(A_j \to A_i) \to \Im(B_j \to B_i) \to C \to 0$
is short exact for all $j \geq i$. Hence the images
$\Im(A_j \to A_i)$ stabilize for $j \geq i$
if and only if the same is true for $\Im(B_j \to B_i)$
as desired.
\end{proof}

\end{document}
